\newcommand\UniversiteAdi{Niğde Ömer Halisdemir Üniversitesi}
\newcommand\BolumAdi{MEKATRONİK BÖLÜMÜ}
\newcommand\DersKodu{MKT2002}
\newcommand\DersAdi{BİLGİSAYARLI KONTROL SİSTEMLERİ}
\newcommand\SinavAdi{Genel Sınav}
\newcommand\SinavTarihi{10.05.2025}
\newcommand\SinavSaati{10:00}
\newcommand\SinavSuresi{90dk}

\pagestyle{fancy}
\fancyhf{} % clear existing header/footer entries
\fancyhead[R]{Öğrenci No:\hspace{4.5cm}}
\fancyhead[L]{Ad Soyad:\hspace{7cm}}
\noindent \includegraphics[width=0.1\textwidth]{logo}
\begin{tabular}{
    p{0.15\linewidth}
    p{0.15\linewidth}
    p{0.2\linewidth}
    p{0.1\linewidth}
    p{0.15\linewidth}}
    \multicolumn{5}{c}{\textbf{\BolumAdi}}\\
    \multicolumn{5}{c}{\textbf{\DersAdi}}\\\hline
    \multicolumn{1}{|r|}{Ders Kodu:}&
    \multicolumn{1}{|c|}{\DersKodu}&
    \multicolumn{1}{|c|}{}& 
    \multicolumn{1}{|r|}{Tarih:}&
    \multicolumn{1}{|c|}{\SinavTarihi} \\\hline
    \multicolumn{1}{|r|}{Sınav Türü:}&
    \multicolumn{1}{|c|}{\SinavAdi}&  
    \multicolumn{1}{|c|}{}&
    \multicolumn{1}{|r|}{Saat:}&
    \multicolumn{1}{|c|}{\SinavSaati}\\\hline
    \multicolumn{1}{|r|}{Dönemi:}&
    \multicolumn{1}{|c|}{2024-2025}&
    \multicolumn{1}{|c|}{}&
    \multicolumn{1}{|r|}{Süre:}&
    \multicolumn{1}{|c|}{\SinavSuresi} \\\hline
    &&&&\\
\end{tabular}\\\\
\noindent\begin{center}
\begin{tabular}{|r|c|c|c|}\hline
    \textbf{Soru:}&
    \textbf{1}&
    \textbf{2}&
    \textbf{Toplam}\\\hline
    \textbf{Puan:}&
    \textbf{50}&
    \textbf{50}&
    \textbf{100}\\\hline
    \textbf{Not:}&50&50&100\\\hline
\end{tabular}\end{center}
\noindent\textbf{Uyarı:}
\begin{itemize}\bfseries
    \item Soruları dikkatlice okuyunuz. Hesap makinesi kullanılabilir.
    \item Defter, kitap ve notlar açık bir sınavdır.
    \item İşlemleri atlamadan ve ayrıntılı olarak veriniz. Sadece nümerik yanıtlar veya çizimler ara işlemler olmadan kabul edilmemektedir.
    \item Yuvarlamalar 2 hane yapılacaktır. $\mathbf{1.99456\approx1.99}$ olarak alınacaktır.
\end{itemize}
    
\begin{enumerate}[\bfseries S1.]
    \item (50p)\,Bir transfer fonksiyonu
    \begin{equation}
        G(s)=\frac{1}{s^2+4s+3}
    \end{equation}
    olarak verilmiştir. Aşımı $\%16.3$ yapacak P kontrolörü tasarlayınız. Aşım kriteri ile
    \begin{equation}
    \begin{split}
        \zeta&=-\frac{\log(os)}{\sqrt{\pi^2+\log(os)^2}}\\
        &=-\frac{\log(0.163)}{\sqrt{\pi^2+\log(0.163)^2}}\\
        &=\frac{1.814}{\sqrt{3.14^2+1.814^2}}\\
        &=\frac{1.814}{\sqrt{9.8696+3.2906}}\\
        &=\frac{1.814}{3.6277}\\
        &=0.5
    \end{split}
    \end{equation}
    ve isterleri sağlayacak polinom
    \begin{equation}
    \begin{split}
        p(s)&=s^2+2\zeta \omega_n s+\omega_n^2\\
        &=s^2+\omega_n s+\omega_n^2
    \end{split}
    \end{equation}
    olarak elde edilmektedir. Kontrolör ile kapalı çevrim
    \begin{equation}
    \begin{split}
        T(s)&=\frac{kG(s)}{1+kG(s)}\\
        &=\frac{k\frac{1}{s^2+4s+3}}{1+k\frac{1}{s^2+4s+3}}\\
        &=\frac{k}{s^2+4s+3+k}
    \end{split}
    \end{equation}
    şeklindedir. Tasarım problemi 
    \begin{equation}
    \begin{split}
        4&=\omega_n\\
        3+k&=\omega_n^2
    \end{split}
    \end{equation}
    olmak üzere çözüm 
    \begin{equation}
    \begin{split}
        \omega_n&=4\\
        k&=13
    \end{split}
    \end{equation}
    şeklinde hesaplanmaktadır. Kapalı çevrim transfer fonksiyonu 
    \begin{equation}
    T(s)=\frac{13}{s^2 + 4 s + 16}
    \end{equation}
    şeklindedir.

    \item (50p)\,Bir transfer fonksiyonu
    \begin{equation}
        G(s)=\frac{1}{s+0.2}
    \end{equation}
    olarak verilmiştir. Aşımı $\%16.3$ ve yerleşme zamanını $t_s=2\,s$ yapan PI kontrolör tasarlayınız. Kapalı çevrim transfer fonksiyonu
    \begin{equation}
    \begin{split}
        T(s)&=\frac{F(s)G(s)}{1+F(s)G(s)}\\
        &=\frac{\frac{k_ps+k_i}{s}\frac{1}{s+0.2}}{1+\frac{k_ps+k_i}{s}\frac{1}{s+0.2}}\\
        &=\frac{k_ps+k_i}{s^2+0.2s+k_ps+k_i}\\
        &=\frac{k_ps+k_i}{s^2+(0.2+k_p)s+k_i}
    \end{split}
    \end{equation}
    olarak hesaplanır. $t_s=2$ isterinden ise 
    \begin{equation}
    \begin{split}
        \omega_n&=\frac{4}{t_s\zeta}\\
        &=\frac{4}{2\cdot 0.5}\\
        &=4
    \end{split}
    \end{equation}
    ve dolayısıyla
\begin{equation}
\begin{split}
    p(s)&=s^2+2\zeta \omega_n s+\omega_n^2\\
    &=s^2+4s+16
\end{split}
\end{equation}
elde edilir. Tasarım problemi
\begin{equation}
\begin{split}
    0.2+k_p&=4\\
    k_i&=16
\end{split}
\end{equation}
ve çözümü
\begin{equation}
\begin{split}
    k_p&=3.8\\
    k_i&=16
\end{split}
\end{equation}
şeklindedir. Kapalı çevrim transfer fonksiyonu
\begin{equation}
    T(s)=\frac{4s+16}{s^2+4s+16}
\end{equation}
olarak elde edilir.
\end{enumerate}
